% pyALDIC User Guide — preamble
% Works with pdflatex, xelatex, or lualatex.

\documentclass[11pt,letterpaper]{article}

% ---------------------------------------------------------------- fonts
\usepackage[T1]{fontenc}
\usepackage[utf8]{inputenc}
\usepackage{lmodern}

% amssymb provides \blacktriangle and \blacktriangledown used in
% sections/05-region-of-interest.tex for the dropdown-button glyph.
\usepackage{amssymb}

% textcomp provides \textdegree and \textmu (used in workflow-type,
% viewing-results, strain-processing, troubleshooting). Modern LaTeX
% bundles these with the kernel, but explicit load keeps the guide
% compilable on older distributions (TeX Live < 2020).
\usepackage{textcomp}

% ---------------------------------------------------------------- layout
\usepackage[margin=1in,headheight=14pt]{geometry}
\usepackage{fancyhdr}
\pagestyle{fancy}
\fancyhf{}
\fancyhead[L]{pyALDIC User Guide}
\fancyhead[R]{\thepage}
\renewcommand{\headrulewidth}{0.4pt}

% ---------------------------------------------------------------- graphics
\usepackage{graphicx}
\graphicspath{{figures/}}

% ---------------------------------------------------------------- colour
\usepackage{xcolor}
\definecolor{aldicblue}{RGB}{65,105,168}
\definecolor{aldicgreen}{RGB}{34,150,94}
\definecolor{aldicred}{RGB}{192,69,72}
\definecolor{codebg}{RGB}{245,245,247}

% ---------------------------------------------------------------- hyperlinks
\usepackage[colorlinks=true,linkcolor=aldicblue,urlcolor=aldicblue,
            citecolor=aldicblue]{hyperref}

% ---------------------------------------------------------------- callouts
\usepackage[most]{tcolorbox}

\newtcolorbox{tipbox}{
  colback=aldicgreen!8!white, colframe=aldicgreen, boxrule=0.5pt,
  arc=1pt, left=4pt, right=4pt, top=3pt, bottom=3pt,
  title={\textbf{Tip}}, fonttitle=\bfseries\small}

\newtcolorbox{warnbox}{
  colback=aldicred!8!white, colframe=aldicred, boxrule=0.5pt,
  arc=1pt, left=4pt, right=4pt, top=3pt, bottom=3pt,
  title={\textbf{Warning}}, fonttitle=\bfseries\small}

\newtcolorbox{notebox}{
  colback=aldicblue!8!white, colframe=aldicblue, boxrule=0.5pt,
  arc=1pt, left=4pt, right=4pt, top=3pt, bottom=3pt,
  title={\textbf{Note}}, fonttitle=\bfseries\small}

% ---------------------------------------------------------------- code
\usepackage{listings}
\lstset{
  basicstyle=\ttfamily\small,
  backgroundcolor=\color{codebg},
  frame=single, framerule=0pt,
  breaklines=true, breakindent=0pt,
  xleftmargin=4pt, xrightmargin=4pt,
  aboveskip=6pt, belowskip=6pt,
}

% ---------------------------------------------------------------- misc
\usepackage{enumitem}
\setlist{topsep=3pt,itemsep=2pt,parsep=0pt}
\usepackage{booktabs}

% ---------------------------------------------------------------- meta
\title{\textbf{pyALDIC User Guide} \\[4pt]
  \large Augmented Lagrangian Digital Image Correlation \\[2pt]
  \normalsize Desktop GUI Walkthrough}
\author{pyALDIC Project}
\date{\today}


\begin{document}
\maketitle
\thispagestyle{empty}

\begin{abstract}
\noindent
This guide walks a new user through a complete pyALDIC analysis,
from first launching the GUI to exporting results. Each section maps
directly to a panel or action in the application so you can follow
along with the real interface open.

If you already know pyALDIC and just need a refresher on one step,
jump to the relevant section via the table of contents.
\end{abstract}

\tableofcontents
\newpage

\section{Overview}
\label{sec:overview}

pyALDIC computes full-field displacement and strain from a sequence
of grayscale images of a deforming specimen. The typical workflow,
from left to right through the GUI, is:

\begin{enumerate}
  \item \textbf{Load images} (a folder of frames, first frame = reference).
  \item \textbf{Choose Workflow Type} --- tracking mode, solver, and
    reference-update policy. These choices are made \emph{before}
    drawing a Region of Interest because they determine which frames
    need one.
  \item \textbf{Draw a Region of Interest} on the required frame(s).
  \item \textbf{Set DIC parameters} --- subset size, node spacing,
    search range, mesh refinement.
  \item \textbf{Run the analysis.} Progress, Cancel, and console log
    are on the right panel.
  \item \textbf{View displacement results} live in the canvas.
  \item \textbf{Compute strain} in the Strain Post-Processing window
    (opens automatically when the run finishes).
  \item \textbf{Export} numerical data, images, or animations.
\end{enumerate}

\begin{notebox}
pyALDIC does \textbf{not} require you to pre-compute calibration,
stereo matching, or any 3D reconstruction. It is strictly 2D DIC
(one camera, planar specimen). If the specimen moves out of plane or
the illumination changes significantly between frames, results will
degrade and the GUI will warn you.
\end{notebox}

\subsection{Three-column layout}

The main window has three columns:

\begin{description}[leftmargin=1.2cm,style=nextline]
  \item[Left sidebar]
    Images panel, \emph{Workflow Type} section, \emph{Region of
    Interest} section, \emph{Parameters} section, and a collapsible
    \emph{Advanced} section.
  \item[Centre canvas]
    Zoom / pan viewport showing the current frame with overlays
    (ROI, mesh, displacement field).
  \item[Right sidebar]
    Run button, progress bar, field selector (disp\_u, disp\_v,
    etc.), visualization controls (colormap, range, opacity),
    physical units, and console log.
\end{description}

\subsection{Typical session times}

On a modern laptop (8-core CPU, 16 GB RAM):

\begin{center}
\begin{tabular}{lll}
  \toprule
  Image size & Grid spacing & Time per frame \\
  \midrule
  $256 \times 256$  & step 8   & $\sim$ 0.2 s \\
  $512 \times 512$  & step 8   & $\sim$ 1 s   \\
  $1024 \times 1024$ & step 4  & $\sim$ 3 s   \\
  \bottomrule
\end{tabular}
\end{center}

Times scale linearly with the number of mesh nodes. The
\emph{Local DIC} solver is roughly $3\times$ faster than
\emph{AL-DIC}.

\section{Installing and launching}
\label{sec:launching}

\subsection{Installation}

Three ways to install, ordered by convenience. Requires
Python 3.10 or newer.

\subsubsection*{Option A: From PyPI (recommended)}

\begin{lstlisting}[language=bash]
pip install al-dic
\end{lstlisting}

This pulls the latest released wheel and every runtime dependency
in one step.

\subsubsection*{Option B: From a GitHub Release wheel}

Useful when you cannot reach PyPI (corporate firewall, air-gapped
network) or when you want to install a specific past version whose
artefacts are still attached to its release page.

\begin{enumerate}
  \item Go to \url{https://github.com/zachtong/pyALDIC/releases}.
  \item Expand \emph{Assets} on the release you want; download the
    \texttt{al\_dic-<version>-py3-none-any.whl} file.
  \item Install locally:
\begin{lstlisting}[language=bash]
pip install ./al_dic-<version>-py3-none-any.whl
\end{lstlisting}
\end{enumerate}

Dependencies are still fetched from PyPI at install time, so this
method only bypasses PyPI for the \texttt{al-dic} package itself.

\subsubsection*{Option C: From source (development)}

Editable install that lets you modify the code and see the effect
without re-installing.

\begin{lstlisting}[language=bash]
git clone https://github.com/zachtong/pyALDIC.git
cd pyALDIC
pip install -e ".[dev]"
\end{lstlisting}

The \texttt{.[dev]} extras include \texttt{pytest} and
\texttt{pytest-xdist} for running the test suite.

\subsection{Launching the GUI}

\begin{lstlisting}[language=bash]
al-dic
# or equivalently
python -m al_dic
\end{lstlisting}

The window opens at a fixed 1380$\times$800 minimum size. You can
resize it to taste; the three columns keep their proportions.

\subsection{What you should see first}

The first time the window opens:
\begin{itemize}
  \item The left image list is empty and shows a drop zone that reads
    \emph{``Drop image folder or Browse''}.
  \item The centre canvas shows an empty dark background.
  \item The right sidebar shows a disabled \textbf{Run DIC Analysis}
    button (you cannot run without images and a Region of Interest).
\end{itemize}

\begin{tipbox}
All user-visible text uses the full phrase \emph{Region of Interest}
instead of the DIC jargon ``ROI''. The 50-pixel Region column in the
image list is abbreviated to ``Region'' to fit, with
Need/Edit/Add button labels inside.
\end{tipbox}

\section{Loading images}
\label{sec:loading-images}

\subsection{Dropping or browsing a folder}

Click the drop zone or drag a folder of image files onto it. Every
image in the folder is loaded in natural sort order. The first image
is always treated as the reference frame (frame 1 in the UI,
index 0 internally).

Supported formats: \texttt{.tif}, \texttt{.tiff}, \texttt{.png},
\texttt{.bmp}, \texttt{.jpg}, \texttt{.jpeg}. Grayscale is assumed;
RGB inputs are converted to luminance on load.

\begin{warnbox}
All images in the folder must have the same dimensions. If they
don't, loading stops and an error dialog tells you which image is
the first mismatch.
\end{warnbox}

\subsection{The image list}

After loading, the left sidebar shows three columns:

\begin{center}
\begin{tabular}{lll}
  \toprule
  Column & Width & Content \\
  \midrule
  \# & fixed & 1-based frame number \\
  Filename & flex & file name without extension \\
  Region & 50 px & Need / Edit / Add button \\
  \bottomrule
\end{tabular}
\end{center}

The \emph{Region} column button changes its label based on state:

\begin{description}[leftmargin=1.2cm,style=nextline]
  \item[\textcolor{aldicred}{\textbf{Need}}]
    This frame is a reference frame (per the current workflow type)
    and does not yet have a Region of Interest.
  \item[\textcolor{aldicgreen}{\textbf{Edit}}]
    This frame already has a Region of Interest; click to modify it.
  \item[Add]
    This frame is \emph{not} a required reference frame, but you can
    optionally add one.
\end{description}

Clicking any row's button activates Region of Interest editing for
that specific frame; the canvas switches to show that frame's
reference and the \emph{EDITING REGION OF INTEREST FOR FRAME xx}
banner appears.

\subsection{Navigating frames}

The frame slider at the bottom of the canvas scrubs through the
sequence. Clicking a row in the image list also jumps to that frame.
Keyboard: left/right arrow keys move by one frame when the canvas
has focus.

\section{Choosing the Workflow Type}
\label{sec:workflow-type}

The \emph{Workflow Type} section is the first choice you make after
loading images. It determines:
\begin{itemize}
  \item How each frame is compared against a reference,
  \item Which solver (Local DIC or AL-DIC) runs,
  \item How often the reference frame refreshes (incremental only),
  \item \textbf{and therefore, which frames need a Region of Interest.}
\end{itemize}

\begin{tipbox}
Get this right before drawing Regions of Interest. The Region of
Interest section below shows a live hint that updates whenever you
change the workflow type, telling you exactly which 1-based frame
numbers require a region.
\end{tipbox}

\subsection{Tracking Mode}

\begin{description}[leftmargin=1.2cm,style=nextline]
  \item[Accumulative]
    Every frame is compared directly against frame 1. Best for small,
    monotonic deformation (strain $<$ 5\%, rotation $<$ 2\textdegree).
    Fast, but fails for large cumulative motion because the reference
    and current subsets become visually too different.
  \item[Incremental]
    Each frame is compared against a rolling reference frame.
    Required for rotations $>$ 5\textdegree{} or large cumulative
    strain. Displacements from consecutive segments are composed into
    cumulative output fields. Slightly slower than accumulative but
    handles arbitrary total deformation.
\end{description}

\subsection{Solver}

\begin{description}[leftmargin=1.2cm,style=nextline]
  \item[Local DIC]
    Each subset is solved independently with IC-GN. Fast, preserves
    sharp local features, but sensitive to noise. Best for high-
    quality images and small deformations.
  \item[AL-DIC]
    Augmented Lagrangian solver that couples local IC-GN with a
    global FEM regularizer. Smoother results, better strain accuracy,
    more robust in noisy or high-gradient regions. $\sim 3\times$
    slower than Local DIC. Default.
\end{description}

When AL-DIC is selected, an indented \emph{ADMM Iterations} spinbox
appears (1--10, default 3). Higher is more regularized but slower.

\subsection{Reference Update (incremental only)}

Visible only when \emph{Tracking Mode} $=$ Incremental. Controls
when the rolling reference frame refreshes:

\begin{description}[leftmargin=1.2cm,style=nextline]
  \item[Every Frame]
    Reference resets every frame. Most robust for large deformation
    (smallest per-step ratio). A Region of Interest is only needed on
    frame 1 --- pyALDIC warps it forward automatically.
  \item[Every N Frames]
    Reference resets every $N$ frames. You need a Region of Interest
    on frame 1, $N+1$, $2N+1$, \ldots. The live hint in the Region of
    Interest section lists the required frames.
  \item[Custom Frames]
    You type a comma-separated list of reference frames, e.g.\
    \texttt{1, 6, 12, 24}. Frame 1 is always included.
\end{description}

\subsection{Which combination to pick?}

\begin{center}
\begin{tabular}{lll}
  \toprule
  Experiment & Tracking & Reference Update \\
  \midrule
  Soft material, $<$ 2\% strain & Accumulative & --- \\
  Soft material, 2--30\% strain & Incremental & Every Frame \\
  Rigid rotation, $>$ 5\textdegree{} & Incremental & Every Frame \\
  Dynamic event / impact     & Incremental & Every Frame \\
  Quasi-static with known    & Incremental & Every N or Custom \\
  \quad load steps           &            & (align $N$ with steps) \\
  \bottomrule
\end{tabular}
\end{center}

\begin{notebox}
When in doubt, use \emph{Incremental + Every Frame}. It is slightly
slower than accumulative but handles every deformation magnitude
well, and only needs one Region of Interest drawn on frame 1.
\end{notebox}

\section{Drawing a Region of Interest}
\label{sec:region-of-interest}

\subsection{Live hint}

The top of the \emph{Region of Interest} section is a blue-accent
info box that updates live as you change the workflow type. It tells
you exactly which frames need a region:

\begin{itemize}
  \item \emph{Accumulative} $\to$ ``Only frame 1 needs a Region of
    Interest. All later frames are compared against it directly.''
  \item \emph{Incremental + Every Frame} $\to$ ``Frame 1 needs a
    Region of Interest. It is automatically warped forward to each
    later frame.''
  \item \emph{Incremental + Every N Frames} $\to$ ``Draw a Region of
    Interest on frames: 1, 6, 11, 16, \ldots''
  \item \emph{Incremental + Custom Frames} $\to$ Lists your custom
    reference frame indices.
\end{itemize}

\subsection{The Region of Interest toolbar}

Three rows of buttons, from top to bottom:

\subsubsection*{Row 1 --- Shape operations}

\begin{description}[leftmargin=1.2cm,style=nextline]
  \item[+ Add $\blacktriangledown$]
    Opens a popup menu: Polygon / Rectangle / Circle. Pick one, then
    draw on the canvas. The drawn region is added (OR'd) to the
    Region of Interest mask.
  \item[\ding{46} Cut $\blacktriangledown$]
    Same menu, but the drawn region is subtracted (AND-NOT) from the
    Region of Interest mask. Use this to cut holes around features
    you don't want to track (bubbles, voids, damaged areas).
  \item[+ Refine $\blacktriangledown$]
    Paint brush for marking mesh-refinement zones (does \emph{not}
    change the Region of Interest --- it marks areas where the
    quadtree mesh should be subdivided for finer resolution).
\end{description}

\begin{warnbox}
The Refine brush is gated to frame 1 only. On any other frame the
button renders with a dashed border and muted colours; hovering shows
the explanation. This is because the brush paints in frame-1 pixel
coordinates and is automatically warped forward --- painting on a
deformed frame would place refinement zones in the wrong material
points.
\end{warnbox}

\subsubsection*{Row 2 --- Import / Batch Import}

\begin{description}[leftmargin=1.2cm,style=nextline]
  \item[Import]
    Load a grayscale / binary image and use it as the Region of
    Interest for the currently-edited frame. Any pixel above zero
    counts as ``inside''.
  \item[Batch Import]
    Opens a dialog to import several mask files at once and assign
    them to specific frames. Useful when you have pre-segmented
    masks from e.g.\ SAM, thresholding, or a dedicated segmentation
    tool.
\end{description}

\subsubsection*{Row 3 --- Save / Invert / Clear}

\begin{description}[leftmargin=1.2cm,style=nextline]
  \item[Save]
    Writes the current frame's Region of Interest mask to a PNG.
  \item[Invert]
    Flips inside/outside on the current frame.
  \item[Clear]
    Empties the current frame's Region of Interest mask. Does not
    delete other frames' regions.
\end{description}

\subsection{Drawing shapes}

After selecting a shape from Add or Cut:

\begin{description}[leftmargin=1.2cm,style=nextline]
  \item[Polygon]
    Click to place vertices, double-click or press Enter to close.
    Right-click to cancel.
  \item[Rectangle]
    Click-drag to define the bounding box.
  \item[Circle]
    Click the centre, drag outward to set the radius, release to
    commit.
\end{description}

While drawing, the toolbar button for the active mode highlights.
The canvas cursor changes to a crosshair. Click \emph{Esc} or the
already-highlighted button again to cancel.

\subsection{Per-frame editing}

Click the Region column button for any frame to jump to editing that
frame's region. The canvas shows the selected frame's image (in
frame-1 coordinates for accumulative mode, or whichever ref frame
you're editing for incremental). A banner overlays the top of the
canvas:

\begin{center}
  \texttt{EDITING REGION OF INTEREST FOR FRAME 07}
\end{center}

Changes are committed as soon as you finish a shape operation.

\section{DIC Parameters}
\label{sec:parameters}

The \emph{Parameters} section of the left sidebar holds the subset-
and mesh-related choices. Change any of these and the downstream
pipeline will recompute when you press Run.

\subsection{Subset Size}

The IC-GN correlation window, in pixels. Displayed as an odd number
(e.g.\ 41), stored internally as the even half-width (40).

\begin{center}
\begin{tabular}{lp{8cm}}
  \toprule
  Value & When to use \\
  \midrule
  21--31 & Fine / dense speckle, small strain. Best spatial resolution. \\
  41 (default) & General speckle patterns with $\sim$ 3--5 px particles. \\
  51--81 & Noisy images, coarse speckle, or low-contrast regions. Lower spatial resolution. \\
  \bottomrule
\end{tabular}
\end{center}

Larger subsets average over more pixels, which improves signal-to-
noise but smooths out local features. Rule of thumb: each subset
should contain 5--10 speckle particles.

\subsection{Subset Step}

Node spacing in pixels, restricted to powers of two:
\texttt{4, 8, 16, 32, 64}. Smaller step means more mesh nodes and a
denser displacement field.

\begin{center}
\begin{tabular}{lp{8cm}}
  \toprule
  Value & Trade-off \\
  \midrule
  4 & Densest grid, very slow. Use only for $\leq$ 512\textsuperscript{2} images. \\
  8 (default for 512\textsuperscript{2}) & Good balance. \\
  16 (default for 1024\textsuperscript{2}) & Standard choice for large images. \\
  32--64 & Very fast, coarse. Only for quick previews. \\
  \bottomrule
\end{tabular}
\end{center}

\begin{tipbox}
Total mesh nodes grow as $(\text{ROI area}) / (\text{step})^2$.
Halving the step quadruples the node count and roughly quadruples
runtime. Start coarse, iterate on Region of Interest and parameters,
then refine the step only when you are happy with the setup.
\end{tipbox}

\subsection{Search Range}

Maximum per-frame displacement (in pixels) that the FFT initial
search can detect. If the real displacement exceeds this, the FFT
returns a clipped peak and pyALDIC auto-expands the search radius
($2\times$ each retry up to image half-size) --- but this costs
time.

\begin{center}
\begin{tabular}{lp{8cm}}
  \toprule
  Value & When to use \\
  \midrule
  8--20 (default 20) & Typical quasi-static experiments. \\
  40--60 & Larger translations, moderate rotations in incremental mode. \\
  80+ & Use only if you know the motion is huge. Diagnose before setting high: a large search range also means more chance of picking wrong NCC peaks in repetitive textures. \\
  \bottomrule
\end{tabular}
\end{center}

\begin{warnbox}
If you see \texttt{FFT search: N nodes have peaks at search boundary}
in the console log, the initial search range is too small. Either
raise it manually or let Auto-expand (in Advanced, on by default)
handle it.
\end{warnbox}

\subsection{Mesh Refinement}

Two independent checkboxes below the search range control quadtree
refinement of the mesh:

\begin{description}[leftmargin=1.2cm,style=nextline]
  \item[Refine Inner Boundary]
    Subdivide mesh elements near \emph{holes} inside the Region of
    Interest --- e.g.\ around a bubble, void, or cut-out region.
  \item[Refine Outer Boundary]
    Subdivide mesh elements near the outer edge of the Region of
    Interest, where boundary effects matter most.
\end{description}

When either box is checked (or the Refine brush has been used on
frame 1), the \emph{Refinement Level} dropdown activates:

\begin{center}
\begin{tabular}{lll}
  \toprule
  Level & Label & Min element size \\
  \midrule
  1 & Light        & $\text{step} / 2$ \\
  2 & Medium       & $\text{step} / 4$ \\
  3 & Heavy        & $\text{step} / 8$ \\
  4 & Extra Heavy  & $\text{step} / 16$ \\
  5 & Ultra        & $\text{step} / 32$ \\
  \bottomrule
\end{tabular}
\end{center}

Available levels depend on subset size and subset step so the min
element size stays reasonable (minimum 2 px element). A live readout
below the dropdown tells you the current minimum element size.

\section{Advanced parameters}
\label{sec:advanced}

The bottom collapsible section on the left sidebar holds parameters
that most users leave at their defaults. Expand it only if you are
debugging convergence or doing cutting-edge experiments.

\subsection{Initial Guess}

Controls how IC-GN gets its starting displacement for each frame.
Four mutually-exclusive options:

\begin{description}[leftmargin=1.2cm,style=nextline]
  \item[Previous frame (default, fastest)]
    Warm-start from the previous frame's converged result. Best for
    smooth, slow motion. Automatically selected when Tracking Mode is
    Accumulative.
  \item[FFT when reference frame updates]
    Run an FFT cross-correlation whenever a new reference frame is
    adopted; warm-start within each segment. This is the default for
    Incremental mode.
  \item[FFT every frame (safest)]
    Full FFT every frame. Eliminates warm-start error propagation
    across frames, but slowest. Use for unpredictable motion
    (oscillation, impact, random).
  \item[Reset FFT every N frames]
    Full FFT every $N$ frames, warm-start in between. Bounds the
    worst-case warm-start error to $N$ frames.
\end{description}

\subsection{Auto-expand search on clipped peaks}

Checkbox, on by default. When the FFT NCC peak lands on the edge of
the search window (indicating the true displacement is larger), the
pipeline automatically retries with $2\times$ search range up to six
times, capped at $\min(H, W) / 2$. Leave this on unless you are
debugging FFT behaviour.

\subsection{Mesh appearance}

Cosmetic only. Choose the colour and line width of the mesh overlay
drawn on the canvas (does not affect computation).

\section{Canvas interactions}
\label{sec:canvas}

The centre of the window is a QGraphicsView-based canvas with
layered overlays. From bottom to top:

\begin{enumerate}
  \item Background image (reference or deformed).
  \item Mesh overlay (when enabled).
  \item Region of Interest outline.
  \item Field overlay (post-run, colour-mapped displacement or
        strain).
  \item Refine-brush overlay (cyan, frame 1 only).
\end{enumerate}

\subsection{Zoom and pan}

\begin{description}[leftmargin=1.2cm,style=nextline]
  \item[Mouse wheel] Zoom in / out at the cursor position.
  \item[Fit] Button at the top of the canvas toolbar --- zoom to fit
    the full image in the viewport.
  \item[100\%] Zoom to 1:1 pixel ratio.
  \item[+ / --] Buttons for discrete zoom steps.
  \item[Middle-click drag] Pan the view.
  \item[Left-click drag in empty area] Pan the view (if no drawing
    tool is active).
\end{description}

\subsection{Frame navigation}

The frame slider at the bottom of the canvas scrubs through the
sequence. The label to the left shows the current frame number and
total count.

A small play/pause button next to the slider auto-advances at a
configurable frame rate --- useful for reviewing displacement fields
after a run.

\subsection{Show on deformed frame}

This toggle lives in the FIELD section of the right sidebar (not
VISUALIZATION, because it controls \emph{where} the field plots,
i.e.\ geometry, rather than how it styles).

\begin{description}[leftmargin=1.2cm,style=nextline]
  \item[Checked (default)]
    The field overlay is plotted on the currently-selected deformed
    frame, with node positions warped by the cumulative
    displacement.
  \item[Unchecked]
    The field overlay is plotted on the reference frame (frame 1)
    with node positions un-warped. Useful for comparing an
    experiment back to the starting configuration.
\end{description}

\section{Running the DIC}
\label{sec:running}

\subsection{The Run button}

The big blue \textbf{Run DIC Analysis} button at the top of the
right sidebar starts the pipeline. It is disabled until:
\begin{itemize}
  \item At least 2 images are loaded,
  \item Frame 1 has a non-empty Region of Interest.
\end{itemize}

If either condition fails, clicking Run (or trying to) raises a
modal dialog telling you exactly what is missing.

\subsection{Cancelling}

Below Run is a single red \textbf{Cancel} button (replaces the older
Pause + Stop pair). It is enabled while the pipeline is running.

Cancelling stops the current run cleanly: already-computed frames
are kept, no further frames are processed, and state returns to
IDLE (not DONE, so the Export button stays disabled).

\subsection{Progress display}

Below the Cancel button:

\begin{itemize}
  \item Thin blue progress bar (0--100 \% of the full pipeline).
  \item Text label with the current stage
        (\emph{Section 3: FFT initial search},
        \emph{Section 6: IC-GN local solve},
        \emph{S8: Frame 5/19}, \ldots).
  \item ELAPSED / REMAINING timestamps.
\end{itemize}

\subsection{Fatal errors}

The console log (bottom of the right sidebar) collects every info,
warning, and error message with a timestamp. \emph{Fatal} errors
(missing images, Region of Interest too small, pipeline-side
exceptions) also trigger a modal dialog. The dialog shows a plain-
English summary; the full traceback stays in the console log.

\begin{warnbox}
The console can scroll off-screen if many info messages accumulate.
Fatal errors will not be missed because of the modal dialog, but
warnings can be. If something looks wrong, check the console log.
\end{warnbox}

\subsection{Completion}

When the run finishes cleanly, the Strain Post-Processing window
opens automatically (see section \ref{sec:strain-processing}). The
Run button re-enables for another run; the Export button becomes
active.

If the run is cancelled or fails, Strain does not open automatically
and Export stays disabled.

\section{Viewing results}
\label{sec:viewing-results}

After a run, the right sidebar populates with field-viewing
controls. The canvas shows a colour-mapped overlay of the selected
field on the background image.

\subsection{FIELD section}

Two toggle buttons (exclusive): \textbf{Disp U} and \textbf{Disp V}.
Selected field is highlighted; clicking the other switches.

Below the toggles, the \emph{Show on deformed frame} checkbox (see
Section~\ref{sec:canvas}) controls whether the overlay plots on the
deformed frame or the reference.

\subsection{VISUALIZATION section}

\subsubsection*{Colormap}

Eight built-in maps: \texttt{jet}, \texttt{viridis}, \texttt{turbo},
\texttt{coolwarm}, \texttt{plasma}, \texttt{inferno},
\texttt{RdBu\_r}, \texttt{seismic}.

\emph{RdBu\_r} and \emph{seismic} are diverging (white at zero);
best for displacement that crosses zero.

pyALDIC remembers the colormap choice per-field. Switching from
\texttt{disp\_u} to \texttt{disp\_v} restores whatever map you set
last time for that field.

\subsubsection*{Color Range}

\begin{description}[leftmargin=1.2cm,style=nextline]
  \item[Auto (default)]
    Range is set from the 2--98 percentile of visible node values
    (excludes extreme outliers but uses real data).
  \item[Manual]
    Uncheck Auto, then type \emph{Min} and \emph{Max} explicitly.
    Useful for reporting (all frames share one range) or when you
    know the physically meaningful bounds.
\end{description}

\subsubsection*{Opacity}

Slider (0--100\%). Default 70\%. Adjusts the alpha blending between
the field overlay and the background image.

\subsection{PHYSICAL UNITS section}

A checkbox plus two fields:

\begin{description}[leftmargin=1.2cm,style=nextline]
  \item[Pixel Size]
    Numeric value, dropdown unit (nm, \textmu m, mm, cm, m, inch).
    Converts raw pixel displacements to physical length on all
    displacement-family fields.
  \item[Frame Rate]
    Used to convert inter-frame velocity ($\Delta U / $ frame) to
    physical speed (length / s).
\end{description}

Strain fields are dimensionless and unaffected by these settings.

\subsection{Console log}

Scrollable read-only widget at the bottom of the right sidebar.
Colour-coded by severity:

\begin{center}
\begin{tabular}{ll}
  \toprule
  Level & Colour \\
  \midrule
  info     & blue / white \\
  warning  & yellow       \\
  error    & red          \\
  success  & green        \\
  \bottomrule
\end{tabular}
\end{center}

Every line is prefixed with a \texttt{HH:MM:SS} timestamp.

\section{Strain post-processing}
\label{sec:strain-processing}

The Strain window is a separate top-level window that opens
automatically when a Run completes. It computes strain from the
stored displacement field, with its own parameter panel and its own
visualization state (independent of the main window).

You can also open it manually via the \textbf{Open Strain Window}
button in the right sidebar.

\subsection{Strain Parameters panel}

Three editors from top to bottom:

\subsubsection*{Method}

\begin{description}[leftmargin=1.2cm,style=nextline]
  \item[Plane fitting (default)]
    Weighted moving least-squares fit of a first-order polynomial to
    the displacement field inside a virtual strain gauge (VSG).
    Gaussian weights with the VSG radius as $\sigma$. Robust,
    smooths local noise.
  \item[FEM nodal]
    Strain from shape-function derivatives of the FEM mesh. Higher
    spatial resolution but more sensitive to per-node noise.
\end{description}

\subsubsection*{VSG size (plane fitting only)}

Virtual Strain Gauge diameter in pixels. Odd integer. Default 41.
Larger = more smoothing. The field reports radius internally as
$(VSG - 1) / 2$.

\subsubsection*{Strain field smoothing}

Gaussian smoothing of the strain tensor field \emph{after}
computation (not a displacement smoother). Dropdown with four
presets and a warning glyph on the aggressive one:

\begin{center}
\begin{tabular}{lll}
  \toprule
  Preset & $\sigma$ relative to step & Behaviour \\
  \midrule
  Off                  & 0    & No smoothing \\
  Light                & 0.5  & Subtle, preserves fine features \\
  Medium               & 1.0  & Balanced (recommended for noisy data) \\
  Strong $\blacktriangle$ & 2.0 & Aggressive, may blur real gradients \\
  \bottomrule
\end{tabular}
\end{center}

Here ``step'' is the DIC node spacing (the value you set as
\emph{Subset Step} in the main window). Smoothing smaller than
$0.5 \times $ step has essentially no effect because the Gaussian
kernel cannot reach any neighbour node; $2 \times $ step is the
upper bound of ``useful'' smoothing.

\subsubsection*{Strain type}

The strain measure the output fields report:

\begin{description}[leftmargin=1.2cm,style=nextline]
  \item[Infinitesimal (default)]
    $\varepsilon = \frac{1}{2}(\nabla u + \nabla u^T)$. Accurate for
    small deformation only. \emph{Not} zero under finite rigid
    rotation --- it reports $\cos\theta - 1$ as false strain.
  \item[Eulerian--Almansi]
    $e = \frac{1}{2}(I - F^{-T} F^{-1})$. Reported in the current
    (deformed) configuration.
  \item[Green--Lagrangian]
    $E = \frac{1}{2}(F^T F - I)$. Reported in the reference
    configuration. \textbf{Exactly zero under rigid rotation.} Use
    this for any experiment with rotation $>$ 2\textdegree{}.
\end{description}

\subsection{Compute Strain button}

Green primary button below the parameters. Click to recompute strain
with the current panel settings. Runs in a background thread; a
progress bar shows per-frame progress.

If you haven't changed any parameter and the strain is already
computed, the stale warning label stays empty. If you change
anything, a yellow \emph{Params changed --- click Compute Strain}
hint appears.

\subsection{Strain field selector}

Below the parameters, a radio-style selector with all available
fields:

\begin{itemize}
  \item \textbf{Displacement} (shared with main window):
        \texttt{disp\_u}, \texttt{disp\_v}, \texttt{disp\_magnitude},
        \texttt{velocity}
  \item \textbf{Strain components}:
        \texttt{strain\_exx}, \texttt{strain\_eyy}, \texttt{strain\_exy}
  \item \textbf{Derived strains}:
        \texttt{strain\_principal\_max},
        \texttt{strain\_principal\_min},
        \texttt{strain\_maxshear},
        \texttt{strain\_von\_mises},
        \texttt{strain\_rotation} (polar-decomposition rotation
        angle in degrees)
\end{itemize}

Visualization controls (colormap, range, opacity) and physical-unit
settings mirror the main window's but are \emph{independent}: the
strain window remembers its own colormap and range choices per
field.

\subsection{Strain vs main window}

\begin{center}
\begin{tabular}{p{4cm}p{4cm}p{4cm}}
  \toprule
  & Main window & Strain window \\
  \midrule
  Shows & Displacement field & Strain fields (plus disp) \\
  State  & Shared AppState & Private viz state \\
  Export button & Export Results & Export Results (same dialog) \\
  Frame navigator & Shared & Private \\
  \bottomrule
\end{tabular}
\end{center}

\section{Exporting}
\label{sec:export}

Both the main window and the Strain window have an \textbf{Export
Results} button. Both open the \emph{same} dialog. The difference is
only which visualization preset (colormap, range, deformed toggle)
seeds the dialog --- whichever window you opened it from.

\subsection{The Export dialog}

Four tabs at the top of the dialog:

\subsubsection*{Data tab}

Check boxes for numeric-data formats:
\begin{itemize}
  \item \textbf{NPZ} --- NumPy compressed archive; preferred for
        Python downstream.
  \item \textbf{MAT} --- MATLAB v7.3 HDF5 file.
  \item \textbf{CSV} --- one row per node per frame; plain text,
        large.
\end{itemize}

Plus a multi-select list of which fields to include (displacement,
strain components, derived strains, rotation).

\subsubsection*{Images tab}

Per-field list with Enable / Colormap / Range / Background options
for rendering each field as PNG / JPG / TIFF / PDF. One file per
frame per enabled field.

\begin{tipbox}
\textbf{Include colorbar} is ON by default. Images without a
colorbar show coloured patches but no scale, which is almost never
useful for quantitative reports. Uncheck it only if you plan to
composite a manually-placed colorbar in Illustrator / Inkscape.
\end{tipbox}

Adjustable settings:
\begin{description}[leftmargin=1.2cm,style=nextline]
  \item[Image format]
    png / jpg / tif / pdf. PNG is lossless and near-universal; use
    PDF for figures you will embed in LaTeX.
  \item[DPI] 72 / 150 / 300 / 600. 300 is print-quality.
  \item[Background]
    \emph{Reference frame}: all frames rendered on the reference.
    \emph{Current frame}: each frame rendered on its deformed image.
\end{description}

\subsubsection*{Animations tab}

Render an MP4 / GIF for each selected field, sweeping through
frames. Configurable FPS (usually match the frame rate of the
experiment). Colorbar on by default here too.

\begin{warnbox}
MP4 export requires \texttt{ffmpeg} on your system PATH. If ffmpeg is
missing, the dialog will switch silently to GIF and log a warning.
\end{warnbox}

\subsubsection*{Report tab}

Generates a multi-page PDF summarizing the run: parameters, key
fields at mid and last frame, RMSE (if ground truth is available),
and optional thumbnail animation.

\subsection{Destination folder}

At the top of the dialog. Defaults to
\texttt{<image\_folder>/al-dic-export/<timestamp>/}. Click
\emph{Browse} to pick a different location. Exports are written
inside that folder with a user-supplied prefix (default
\texttt{result}).

\subsection{Running an export}

Click the primary \textbf{Export} button at the bottom. Progress
shows per-file (data) and per-frame (images / animations). The
Cancel button stops cleanly; already-written files are kept.

When the export completes, a green success message tells you how
many files were written and their location. The dialog stays open
so you can export more formats; click Close when done.

\section{Sessions --- save and resume}
\label{sec:session}

pyALDIC can serialize your entire setup (image folder reference,
parameters, Regions of Interest, physical units) into a single
JSON file so you can close the application and pick up exactly
where you left off.

\subsection{File format}

\begin{itemize}
  \item Extension: \texttt{.aldic.json}
  \item Single-file portable: the Regions of Interest are stored
        inline as base64-encoded PNG.
  \item Versioned: the top-level key \texttt{schema\_version}
        protects against silent misinterpretation of newer schemas.
\end{itemize}

\subsection{What is saved}

\begin{itemize}
  \item Image folder path and the ordered list of filenames.
  \item Every per-frame Region of Interest mask.
  \item Core DIC parameters: subset size, step, search range,
        tracking mode, reference-update policy, solver choice,
        refinement settings, initial-guess mode.
  \item Physical units (pixel size, unit, frame rate).
\end{itemize}

\subsection{What is not saved}

\begin{itemize}
  \item Pipeline results (re-run with Run --- cheap vs.\ re-
        configuring).
  \item Run state / progress / elapsed time.
  \item Transient visualization preferences (colormaps, opacity,
        colour ranges).
\end{itemize}

\subsection{Menu}

\begin{description}[leftmargin=1.2cm,style=nextline]
  \item[File $\to$ Save Session\ldots] (\texttt{Ctrl+S})
    Opens a file dialog. Suggested extension
    \texttt{*.aldic.json}. Overwrites if the file exists.
  \item[File $\to$ Open Session\ldots] (\texttt{Ctrl+O})
    Opens a file dialog. Filter: \texttt{pyALDIC Session
    (*.aldic.json)}.
\end{description}

\subsection{Opening a session whose images moved}

If the image folder path in the session no longer exists, pyALDIC
does \emph{not} fail the load. Instead it:
\begin{enumerate}
  \item Restores parameters and Regions of Interest.
  \item Logs a warning: ``Image folder ... no longer exists;
        parameters and Regions of Interest were restored but images
        must be re-selected manually.''
  \item Leaves the image list empty.
\end{enumerate}

You can then drop a new folder of the same images and the Regions
of Interest will still be aligned.

\begin{tipbox}
This behaviour is deliberate so you can share a \texttt{.aldic.json}
file with a colleague whose images live at a different path.
They re-drop the images; all parameters and Regions of Interest
come along.
\end{tipbox}

\subsection{Corrupt or future-version files}

Malformed JSON, non-object root, or an unknown
\texttt{schema\_version} raise a modal error dialog. The fatal-
error path is used so you cannot miss the failure. The application
keeps whatever state it had before the attempted load; nothing is
silently replaced.

\section{Troubleshooting}
\label{sec:troubleshooting}

\subsection{``Cannot start analysis''}

A modal dialog with one of these messages:

\begin{description}[leftmargin=1.2cm,style=nextline]
  \item[Need at least 2 images]
    Drop an image folder into the Images panel.
  \item[Define a Region of Interest first]
    Use the Region of Interest toolbar to draw or import one on
    frame 1.
  \item[Region of Interest is empty]
    You drew something, then clicked Cut over the whole region.
    Click Clear or redraw.
  \item[Region of Interest too small]
    The bounding box is smaller than $2 \times \text{Subset Step}$.
    Enlarge the Region or reduce Subset Step.
\end{description}

\subsection{``DIC analysis failed''}

Shows an exception type and message. The full traceback stays in
the console log --- scroll up to find it. Common causes:
\begin{itemize}
  \item Images of different sizes (should be caught at load; if you
        see it at Run, report a bug).
  \item Extremely low subset correlation (black-on-black image or
        no speckle). Check your input data first.
\end{itemize}

\subsection{Slow run}

\begin{itemize}
  \item Subset Step too small: halve the step $\to$ 4$\times$ node
        count $\to$ 4$\times$ runtime. Start with step = 16 for
        large images.
  \item AL-DIC solver: $\sim 3\times$ slower than Local DIC.
        Preview with Local DIC, switch to AL-DIC for the final run.
  \item Search Range much larger than actual motion: every FFT call
        costs proportional to (search)$^2$. Set it just above your
        expected inter-frame motion.
\end{itemize}

\subsection{Poor results in specific regions}

\begin{itemize}
  \item Noisy or fuzzy edges of the Region of Interest: enable
        \emph{Refine Outer Boundary} to densify mesh near the edge.
  \item Hole boundaries (e.g.\ bubble edges): enable \emph{Refine
        Inner Boundary}.
  \item Whole regions that simply failed to converge: check the
        console log for ``n\_bad\_pts'' warnings. Usually means
        Search Range was too small or the speckle was washed out
        there.
\end{itemize}

\subsection{Large rotation gives crazy strain values}

You are probably using \emph{Infinitesimal} strain with a rotation
$>$ 2\textdegree. Infinitesimal strain interprets rigid rotation as
strain of magnitude $\cos\theta - 1$ (about $-0.13$ at 30\textdegree).
Switch the Strain window's \emph{Strain type} to
\textbf{Green--Lagrangian}, which is exactly zero under any rigid
rotation.

\subsection{Refine brush grayed out}

Brush is gated to frame 1 only. Click into the image list or use
the frame slider to return to frame 1. The button re-enables
automatically.

\subsection{Session file won't open}

The fatal-error dialog shows the specific reason:
\begin{itemize}
  \item \texttt{not valid JSON}: file corrupted or hand-edited
        wrong. Open it in a text editor, validate against a JSON
        linter.
  \item \texttt{Unsupported schema version}: the file was saved by a
        newer pyALDIC version than the one you are running. Upgrade
        pyALDIC.
  \item \texttt{Image folder ... no longer exists}: this is a
        warning, not a failure. Parameters and Regions of Interest
        still load; re-select the image folder manually.
\end{itemize}


\end{document}
